% ============================================================
% chapter-generalities.tex
% Capítulo 1: Generalidades
% ============================================================

\chapter[Generalidades]{Generalidades de lassssssssssssssssssssssssssssssss
ssssssssssssssssssssssssss
sssssssssssssssssssssssssssssssssss}
\label{ch:generalities}

\section{Problemática}

\subsection{Árbol de Problemas}
[Incluir diagrama del árbol de problemas.]
\lipsum

\subsection{Descripción}
[Describir el problema. Sustentar afirmaciones con referencias \citex.]

\subsection{Problema seleccionado}
[Enunciar el problema seleccionado.]

\section{Objetivos}

\subsection{Objetivo general}
[Enunciar el objetivo general.]

\subsection{Objetivos específicos}
\begin{enumerate}
  \item[O1.] [Descripción del objetivo 1.]
  \item[O2.] [Descripción del objetivo 2.]
  \item[O3.] [Descripción del objetivo 3.]
\end{enumerate}

\subsection{Resultados esperados}
\begin{enumerate}
  \item[O1.] [Descripción del objetivo 1.]
  \begin{itemize}
    \item [Resultado 1 del objetivo 1.]
    \item [Resultado 2 del objetivo 1.]
  \end{itemize}
  \item[O2.] [Descripción del objetivo 2.]
\end{enumerate}

\subsection{Mapeo de objetivos, resultados y verificación}

\begin{table}[H]
  \centering
  \caption{Mapeo de objetivos, resultados y verificación}
  \label{tab:mapping}
  \begin{tabular}{|p{4cm}|p{4cm}|p{5cm}|}
    \hline
    \textbf{Resultado} & \textbf{Medio de verificación} & \textbf{Indicador objetivamente verificable} \\
    \hline
    [Resultado 1] & [Medio] & [Indicador] \\
    \hline
  \end{tabular}
\end{table}

\section{Métodos y Procedimientos}
[Describir para cada resultado la herramienta o método a usar,
con referencias donde se encuentren más detalles \citex.]
